% Independent convention audit; the manuscripts are not modified by this file.
\section*{BPZ and plumbing with the same inversion}

At a first vertex slot use the local BPZ map and its spin-frame lift
\begin{equation}
 I:z\longmapsto1/z,\qquad \lambda_I(z)=-i/z.
 \label{eq:new-bpz-map}
\end{equation}
At an edge use
\begin{equation}
 I_q:z\longmapsto q/z,\qquad \lambda_{I_q}(z)=-i\sqrt q/z,
 \qquad \sqrt q\sim1\quad(q\sim1).
 \label{eq:new-edge-map}
\end{equation}
The choice $\sqrt{-1}=i$ makes $-i/z$ a square root of
$d(1/z)/dz=-1/z^2$.  The edge map is the BPZ map followed by the
spin-frame scaling $(w,\vartheta)\mapsto(qw,\sqrt q\,\vartheta)$.
At an NS puncture $\lambda$ is the multiplier of the odd coordinate;
at a Ramond puncture it specifies the spin-frame transport on the
annulus and does not assert the same formula for a Ramond-adapted odd
coordinate at the divisor.
Consequently the edge propagator is $q^{L_0}$, with the specified lift
sign $\eta_e^{F}$ if requested.  Removing the primary factor $q^{h_e}$
leaves $q^{|A|}$.  There is no additional half-turn at an edge.

The geometric BPZ transpose is graded and linear:
\begin{equation}
 (AB)^\dagger=(-1)^{|A||B|}B^\dagger A^\dagger,
 \qquad
 L_n^\dagger=L_{-n},\quad
 G_r^\dagger=iG_{-r},\quad
 \psi_r^\dagger=-i\psi_{-r}.
 \label{eq:new-bpz-modes}
\end{equation}
In particular $L_n^\dagger=L_n$ only for $n=0$.  The factors on
the three currents follow from their weights $2,3/2,1/2$ and the
chosen spin lift: contour reversal under $z\mapsto1/z$ cancels the
Jacobian reversal, leaving $(-i)^{2h}A_{-r}$ for a field of weight
$h$.  They preserve the SCA and Clifford relations under
the graded transpose; applying the transpose twice gives $(-1)^F$.
Our bilinear pairing obeys
\begin{equation}
 \langle Ax,y\rangle_{\rm BPZ}
 =(-1)^{|A||x|}\langle x,A^\dagger y\rangle_{\rm BPZ}.
 \label{eq:new-bpz-pairing}
\end{equation}
Set the even Ramond ground norm to one.  The paper's
$G_0w^\pm=i\beta e^{\mp i\pi/4}w^\mp$ and the symmetric auxiliary
$\psi_0u^a=u^{1-a}/\sqrt2$ then imply
\begin{equation}
 \langle w^\alpha,w^\gamma\rangle_{\rm BPZ}
 =\operatorname{diag}(1,-1)_{\alpha\gamma},\qquad
 \langle u^a,u^b\rangle_{\rm BPZ}
 =\operatorname{diag}(1,-i)_{ab}.
 \label{eq:new-ground-gram}
\end{equation}
All off-diagonal ground pairings vanish.  Equation
\eqref{eq:new-bpz-pairing}, not an ungraded reversal of oscillator
words, defines every descendant Gram matrix.  In particular the
inverse Gram matrices in a sewn block are the inverses of these
geometric matrices.  For the paper's PBW word
$\mathbf L_{-A}=L_{-n_1}\cdots L_{-n_\ell}
G_{-r_1}\cdots G_{-r_p}$ and a ground state $g$, the rule is
\begin{equation}
 \langle\mathbf L_{-A}g,y\rangle_{\rm BPZ}
 =i^p(-1)^{p|g|+p(p-1)/2}
  \langle g,G_{r_p}\cdots G_{r_1}
                   L_{n_\ell}\cdots L_{n_1}y\rangle_{\rm BPZ}.
 \label{eq:new-pbw-gram}
\end{equation}
The rightmost pairing uses the ground metric above.  Direct evaluation
of \eqref{eq:new-pbw-gram} fixes the open R-edge tensor as well as
closed blocks.  For the even NS primaries used in the paper,
the ground values of the ordered three-point forms remain the
paper's values.  Their descendants use the same BPZ map as the Gram
matrix.  Explicitly, with insertions at $(\infty,1,0)$, we fix
\begin{align}
 &\rho_0(\phi_1,\phi_2,\phi_3)
 =\rho_1(\phi_1,G_{-1/2}\phi_2,\phi_3)=1,\\
 &\rho^{(\eta)}_0(\phi_1,w^+_2,w^+_3)=1,\qquad
 \rho^{(\eta)}_0(\phi_1,w^-_2,w^-_3)=\eta,\\
 &\rho^{(\eta)}_1(\phi_1,w^+_2,w^-_3)=1,\qquad
 \rho^{(\eta)}_1(\phi_1,w^-_2,w^+_3)=i\eta,\\
 &\rho_{\mathsf F}(\mathbf1,\mathbf1,\mathbf1)
 =\rho_{\mathsf F}(\mathbf1,u^0,u^0)=1,\qquad
 \rho_{\mathsf F}(\mathbf1,u^1,u^1)=i.
 \label{eq:new-ground-vertices}
\end{align}
These choices fix the ordered forms; the spin lift and Ward
identities below fix all their descendant values.  In particular,
the local form at infinity uses $I$ itself, without a separate
conversion at a sewn edge.
If the NS primary in the first slot is assigned odd parity, the
geometric BPZ norm of that ground state is $i$ and each displayed
first-slot ground value acquires $i$.  This choice makes the
graded transpose \eqref{eq:new-bpz-pairing} hold for all descendants;
it also prevents an odd primary from being silently treated as even
in a PBW comparison.

For clarity, write $p_i$ for the parity of the state in slot $i$.
The ordered NS--R--R supercurrent Ward identity in this convention is
\begin{align}
 &\sum_{j\ge0}(-1)^j\binom{n+\frac12}{j}
 \rho_f^{(\eta)}(G_{j-m-n-\frac12}x_1,x_2,x_3)\nonumber\\
 &=i(-1)^{p_1}\sum_{j\ge0}\binom{m+\frac12}{j}
 \rho_f^{(\eta)}(x_1,G_{n+j}x_2,x_3)\nonumber\\
 &\quad-(-1)^{p_1+f+p_2+n}\sum_{j\ge0}(-1)^j
 \binom{n+\frac12}{j}
 \rho_f^{(\eta)}(x_1,x_2,G_{m+j}x_3).
 \label{eq:new-rr-g-ward}
\end{align}
Here the $f$ in the origin term is the parity of the ordered
NS--R--R form; it is needed for its odd ground values.  The displayed
identity assumes the even NS primary of the paper.  If the first-slot
NS primary is odd, reverse the sign of the origin term, while keeping
$p_i$ as the total state parities elsewhere.  This extra sign occurs
because $f$ labels the Ramond ground pairing, not the intrinsic parity
of the NS primary.  With the same
branch $\sqrt{z(z-1)}$ as the paper, the auxiliary identity is
\begin{align}
 &\sum_{j\ge0}(-1)^j\binom{n-\frac12}{j}
 \rho_{\mathsf F}(\psi_{j-m-n+\frac12}a_1,a_2,a_3)\nonumber\\
 &=-i(-1)^{p_1}\sum_{j\ge0}\binom{m-\frac12}{j}
 \rho_{\mathsf F}(a_1,\psi_{n+j}a_2,a_3)\nonumber\\
 &\quad-(-1)^{p_1+p_2+n}\sum_{j\ge0}(-1)^j
 \binom{n-\frac12}{j}
 \rho_{\mathsf F}(a_1,a_2,\psi_{m+j}a_3).
 \label{eq:new-rr-psi-ward}
\end{align}
The parity $p_i$ in \eqref{eq:new-rr-psi-ward} refers to $a_i$.
For all-NS vertices, omit the Ramond square-root multiplier.  The
corresponding two identities are
\begin{align}
 &\sum_{j\ge0}(-1)^j\binom nj
 \rho_a(G_{j-m-n+\frac12}x_1,x_2,x_3)\nonumber\\
 &=i(-1)^{p_1}\sum_{j\ge0}\binom mj
 \rho_a(x_1,G_{n+j-\frac12}x_2,x_3)\nonumber\\
 &\quad+i(-1)^{p_1+p_2+n}\sum_{j\ge0}(-1)^j\binom nj
 \rho_a(x_1,x_2,G_{m+j-\frac12}x_3),
 \label{eq:new-ns-g-ward}\\
 &\sum_{j\ge0}(-1)^j\binom nj
 \rho_{\mathsf F}(\psi_{j-m-n-\frac12}a_1,a_2,a_3)\nonumber\\
 &=-i(-1)^{p_1}\sum_{j\ge0}\binom mj
 \rho_{\mathsf F}(a_1,\psi_{n+j+\frac12}a_2,a_3)\nonumber\\
 &\quad-i(-1)^{p_1+p_2+n}\sum_{j\ge0}(-1)^j\binom nj
 \rho_{\mathsf F}(a_1,a_2,\psi_{m+j+\frac12}a_3).
 \label{eq:new-ns-psi-ward}
\end{align}
The Virasoro Ward identity is unchanged, since
$L_n^\dagger=L_{-n}$ and $L_n$ is even.

With the auxiliary factor ordered first, the BPZ pairing on a tensor
state is the graded product
\begin{equation}
 \langle a\otimes x,b\otimes y\rangle_{\rm BPZ}
 =(-1)^{|a||x|}
   \langle a,b\rangle_{\mathsf F}\langle x,y\rangle_{\mathsf{SCA}}
 \label{eq:new-tensor-gram}
\end{equation}
on the nonzero parity-diagonal part.  This sign makes the tensor
pairing obey \eqref{eq:new-bpz-pairing} for $G$ and $\psi$ acting with
the usual Koszul rule.  The ordered enlarged local form is
\begin{equation}
 \widehat\rho(a_1\otimes x_1,a_2\otimes x_2,a_3\otimes x_3)
 =(-1)^{|a_1||x_1|+|x_2||a_3|}
   \rho_{\mathsf F}(a_1,a_2,a_3)\rho(x_1,x_2,x_3).
 \label{eq:new-enlarged-vertex}
\end{equation}
The first sign is the graded BPZ dual at infinity; the second moves
the middle physical operator past the third auxiliary operator.
Ground values and all descendant values are fixed by the component
Ward identities and these ordering rules.  In the oscillator
implementation, the local change from its branching frame to the
ordered $\rho_f^{(\eta)}$ frame is still made before applying
\eqref{eq:new-enlarged-vertex}; it is a basis change, not a new
physical three-point constant.

For any plumbing graph use exactly the paper's PBW state sum, with
each inverse Gram replaced by the inverse of
\eqref{eq:new-bpz-pairing}, each local form by the BPZ form above,
and the original Koszul permutation sign $(-1)^{K_\Gamma}$.
Writing $\Phi_{v_i}$ for the selected SCA state on a half-edge,
the ordinary block is
\begin{align}
 \mathbf F_{\hat\Gamma}^{\{f,\eta,a\}}
 &=\sum_{\{A,\alpha\}}(-1)^{K_\Gamma(\epsilon)}
 \prod_e q_e^{|A_{e_1}|}\eta_e^{\epsilon_e}
 \prod_{e\in E_{\NS}}\mathbf G_{h_e}^{A_{e_1}A_{e_2}}
 \prod_{e\in E_{\R}}\mathbb G_{\beta_e}^{A_{e_1}\alpha_{e_1},A_{e_2}\alpha_{e_2}}\nonumber\\
 &\quad\times\prod_{v\in V_{\NS}}
       \rho_{a_v}(\Phi_{v_1},\Phi_{v_2},\Phi_{v_3})
 \prod_{v\in V_{\R}}\rho_{f_v}^{(\eta_v)}
       (\Phi_{v_1},\Phi_{v_2},\Phi_{v_3}).
 \label{eq:new-physical-block}
\end{align}
The inverse Gram matrices here use \eqref{eq:new-bpz-pairing};
this incorporates the lift of $q_e/z$ into the state sum.
The permitted states and parities obey the selection rules of the
paper's block definition.
For the enlarged block use the tensor inverse Gram from
\eqref{eq:new-tensor-gram} and the vertices
\eqref{eq:new-enlarged-vertex}.  Explicitly, the definition is the
same sum as \eqref{eq:new-physical-block}, now over both $A$ and
$\mathsf A$, with
\begin{equation}
 (-1)^{K_\Gamma(\epsilon+\epsilon')}
 \prod_e q_e^{|A_{e_1}|+|\mathsf A_{e_1}|}
            \eta_e^{\epsilon_e+\epsilon'_e}
 \prod_e\widehat G_e^{(A,\mathsf A)_{e_1},(A,\mathsf A)_{e_2}}
 \prod_v\widehat\rho_v
 \label{eq:new-enlarged-block}
\end{equation}
in each summand.  Here $\widehat G_e$ is the inverse of the full
graded tensor Gram matrix, and $\widehat\rho_v$ has the three states
specified by the corresponding term of the sum.  This formula
contains no independent $(-1)^{\epsilon'_e}$ sewing factor; the
fermion signs are fixed by the BPZ inverse metric.
On a split edge with a middle
$\Theta\psi$ insertion, apply this rule to the split graph first;
only then take the diagonal in the two middle plumbing powers.
On that middle sphere use
$\Theta u^{\mathsf b}=(-1)^{\mathsf b+1}u^{1-\mathsf b}$, with
$\Theta$ anticommuting with $\psi_r$ and $G_r$ and commuting with
$L_n$.  The normalization of $\psi$ is the state
$v_{1/2}$ at $P=Q/2$.  The local matrix element in the
double-Virasoro basis is
\begin{equation}
 \mathbb S(n_R,n_L,\alpha)
 =\frac{\langle\!\langle v_{n_R}^{\alpha}|\Theta\psi(1)
                |v_{n_L}^{\alpha}\rangle_{\rm BPZ}}
        {\|v_{n_R}^{\alpha}\|_{\rm BPZ}
         \|v_{n_L}^{\alpha}\|_{\rm BPZ}}.
 \label{eq:new-splitting}
\end{equation}
Its descendant values and the recursion in $n_L,n_R$ follow from
the same $L_{\pm1}$ Ward identity, with the new BPZ pairing in the
matrix element.  The double-Virasoro expansion of the split block is
\eqref{eq:new-double-virasoro} on the split graph, with
$\mathbb S$ at every inserted sphere.  Diagonal extraction then
replaces $q_{e_L}^m q_{e_R}^m$ by $q_e^m$ and keeps equal lift
parities.  This is also the order of operations for its convolution
with the inserted free-fermion block.

Expanding the enlarged state sum gives a convolution of the physical
and auxiliary blocks.  If their parities on edge $e$ are $\epsilon_e$
and $\epsilon'_e$, the convolution sign, in addition to the two
individual block signs, is
\begin{equation}
 (-1)^{K_\Gamma(\epsilon+\epsilon')+K_\Gamma(\epsilon)
          +K_\Gamma(\epsilon')
          +\sum_e\epsilon_e\epsilon'_e
          +\sum_v(\epsilon_{v_1}\epsilon'_{v_1}
                   +\epsilon_{v_2}\epsilon'_{v_3})}.
 \label{eq:new-convolution-sign}
\end{equation}
All exponents are modulo two.  This follows term by term from
the Koszul permutation, the tensor inverse Gram, and the enlarged
local form, respectively.  The same formula on the split graph
defines the inserted convolution; it does not require inventing a
plumbing parameter on the unsplit edge.
For the genus-two theta graph, the two first-slot vertex signs in
\eqref{eq:new-convolution-sign} cancel.  In the parity arrays used by
the code, the product of monomials is therefore
\begin{align}
 &\left(\prod_e q_e^{r_e}\eta_e^{\epsilon_e}\right)
 \star\left(\prod_e q_e^{s_e}\eta_e^{\epsilon'_e}\right)\nonumber\\
 &\quad=(-1)^{K_\Gamma(\epsilon+\epsilon')+K_\Gamma(\epsilon)
                +K_\Gamma(\epsilon')+\sum_e\epsilon_e\epsilon'_e}
 \prod_e q_e^{r_e+s_e}\eta_e^{\epsilon_e+\epsilon'_e}.
 \label{eq:new-theta-convolution}
\end{align}
The direct fermion state sum supplies $\mathbf F^{\mathsf F}$,
and the coefficientwise restricted convolution recovery uses this
product.
At zero plumbing order its only theta parity components are $000$
and $110$: the latter is the negative of the former without the
insertion, and equals it with the insertion.  Hence this constant
has no inverse on the full parity algebra.  The code subtracts
already determined positive-level products, divides by twice the
$000$ coefficient on the recoverable sector, and checks the sector
relation at each coefficient.  The odd-loop insertion changes which
sector is recoverable; the physical block remains the same object.

The algebraic $\mathsf{Vir}\oplus\mathsf{Vir}$ generators and their
central charges are unchanged.  Their modes satisfy
$(L_n^{(j)})^\dagger=L_{-n}^{(j)}$ under the graded pairing, since
the $i$ and $-i$ in $G^\dagger$ and $\psi^\dagger$ cancel with graded
reversal.  The complete double-Virasoro sum has the paper's form,
\begin{align}
 \widehat{\mathbf F}_{\hat\Gamma}^{\{f,\eta,a\}}
 &=\sum_{\{n,\alpha\}}(-1)^{K_\Gamma(\epsilon)}
   \prod_{e\in E_{\NS}}q_e^{2n_e^2}\eta_e^{2n_e}
   \prod_{e\in E_{\R}}q_e^{2n_e^2-1/8}\eta_e^{\alpha_e}
   \prod_{j=1}^2 F_\Gamma(h_{e,n_e}^{(j)},c^{(j)})\nonumber\\
 &\quad\times\prod_{v\in V_{\NS}}\mathbf B_{a_v}(n_{v_1},n_{v_2},n_{v_3})
   \prod_{v\in V_{\R}}\mathbb B_{f_v}^{(\eta_v)}
   (n_{v_1},n_{v_2},n_{v_3},\alpha_{v_2},\alpha_{v_3}).
 \label{eq:new-double-virasoro}
\end{align}
The branching coefficients are the enlarged three-point forms of
$v_n$ and $v_n^\alpha$, divided by the corresponding BPZ norm
roots.  The $L_{\pm1}$ branching recursion still follows from the
Virasoro Ward identity because $L_n^\dagger=L_{-n}$.
Thus the CCY Virasoro recursion and Schottky large-$c$
seed are unchanged.  The branching norms are evaluated with
\eqref{eq:new-tensor-gram}, and a branching vertex uses
\eqref{eq:new-enlarged-vertex}.  Sew the resulting Virasoro blocks
with the new inverse norms, divide by the auxiliary block using
\eqref{eq:new-convolution-sign}, and retain the indicated restricted
sector when an odd Ramond loop makes the ordinary auxiliary factor
noninvertible.

\paragraph{Directed evidence.}
The independent local residue test
\path{C++/tests/bpz_one_over_z_ward.cpp} verifies the new
NS--R--R $G$ and $\psi$ identities at 40 digits.  The BPZ-native
physical Gram and vertex path in \path{C++/include/ramond/direct_pbw.hpp}
is compared with the double-Virasoro output by
\path{C++/tests/partition_blocks.cpp} through total level three,
including both parities of the NS primary and the inserted odd-loop
sector.  The native convolution is checked coefficientwise there.
The script \path{C++/tools/check_theta_pbw_oracle.py} compares against
the separate theta implementation of the SCblock notes; that oracle
is not Yutai's sphere or torus calculation.  The actual Type 0B
sphere four-point block arrays are tested by
\path{C++/tools/check_yutai_sphere_blocks.py}, after putting their
ordered three-point forms in a common frame.  For Yutai's torus
two-point calculation, \path{C++/tools/check_yutai_torus_open_bpz.py}
compares the actual open-R-edge block coefficients, with two sewn
edges, through NS twice-level two and R level one.  The external
punctures include ground states and one $G_{-1}$ insertion on either
side.  All 576 comparisons agree to $2.78\times10^{-16}$ after the
BPZ phase $i^{\epsilon_{\NS}-\epsilon_{\R}}$ from the two sewn edges;
here $\epsilon_{\R}$ is the parity of the complete propagating R
state.  The raw
arrays differ by as much as $\sqrt2$.  This is a block check, not a
recalculation of a full torus amplitude.
Fixed-spin genus-two partition comparisons are finite-cutoff numerical
checks, not analytic proofs of exact cross-channel equality.  A fresh
native-BPZ run at 40 digits on ten matched surfaces used total level
five for the NS--R--R source, level eight for the all-NS target,
$7^3$ and $10^3$ momentum nodes respectively, and two individually
transported spins on each surface.  Its twenty source/target ratios
lie between $0.9998405584$ and $1.0000233972$; the maximum deviation
from one is $1.59442\times10^{-4}$, and the wall time was $732.6$ s.
The inputs have finite quadrature, geometry, and structure-constant
accuracy, so the residual is not solely a conformal-block error.
